%------------------------------------------------------------------------------%
% FREDHOLM OPERATORS AND THE ESSENTIAL SPECTRUM                                %
%------------------------------------------------------------------------------%
% File  : Fredholm_Operators_and_the_Essential_Spectrum.tex                    %
% Author: Klaus Hoeflinger, Beethovenstrasse 11, D-73117 Wangen                %
% Email : khoeflinger@t-online.de                                              %
%------------------------------------------------------------------------------%

%------------------------------------------------------------------------------%
%                       DOCUMENT CLASS and INCLUDE FILES                       %
%------------------------------------------------------------------------------%
\documentclass[11pt,twoside]{article}
\usepackage{geometry}
\geometry{paperheight=241mm,
          paperwidth=152mm,
          top=22mm,
          bottom=17mm,
          left=20mm,
          right=20mm,
          textwidth=112mm,
          textheight=202mm,
          headheight=10mm,
          headsep=3mm,
          footskip=9mm,
          includehead,
          includefoot,
          marginparsep=0mm,
          marginparwidth=0mm}

\renewcommand{\baselinestretch}{0.945}\normalsize

\AtBeginDocument{\setlength\abovedisplayskip{2mm}}
\AtBeginDocument{\setlength\belowdisplayskip{2mm}}

%------------------------------------------------------------------------------%
% define title formats                                                         %
%------------------------------------------------------------------------------%
\usepackage{titlesec}

\renewcommand{\thesection}{\arabic{section}}
\renewcommand{\sfdefault}{FiraSans}
\renewcommand{\ttdefault}{pcr}
\renewcommand{\rmdefault}{antiqua}

\titleformat{\part}[display]
{\normalfont \large \rmfamily \bfseries \centering}
{\small{\thepart.}}{2pt}{}

\titleformat{\section}[runin]
{\normalfont \rmfamily \bfseries}
{\thesection}{1pt}{.\;}[.]

%\titleformat{\section}
%{\normalfont \bfseries \small \filcenter}
%{\thesection.\;}{0mm}{}

\titleformat{\subsection}[runin]
{\normalfont \itshape}
{}{0mm}{}

\titleformat{\subsubsection}[runin]
{\normalfont \itshape}
{}{}{}{}

\titleformat{\paragraph}[runin]
{\normalfont}
{}{}{}{}

\titleformat{\subparagraph}[runin]
{\normalfont}
{}{}{}{}

\titlespacing*{\part}          { 0pt}{ 0pt}{ 8pt}
\titlespacing*{\section}       { 0pt}{ 7pt}{ 4pt}
\titlespacing*{\subsection}    { 0pt}{14pt}{ 6pt}
\titlespacing*{\subsubsection} { 0pt}{ 6pt}{12pt}
\titlespacing*{\paragraph}     { 0pt}{ 6pt}{12pt}
\titlespacing*{\subparagraph}  { 0pt}{ 0pt}{12pt}

%------------------------------------------------------------------------------%
% define enumerate parameters                                                  %
%------------------------------------------------------------------------------%
\usepackage{enumitem}
\setlength{\parindent}{3pt}
\setlength{\parskip}  {0pt}

%\setlist{noitemsep}

\setlist[enumerate]{
         leftmargin=12pt,
         rightmargin=0pt,
         itemsep=1pt,
         itemindent=3pt,
         topsep=3pt,
         labelsep=4pt,
         partopsep=0pt,
         labelwidth=0pt,
         listparindent=0pt,
         parsep=0pt}

\setlist[itemize]{
         leftmargin=12pt,
         rightmargin=0pt,
         itemsep=0pt,
         itemindent=1pt,
         topsep=1pt,
         labelsep=4pt,
         partopsep=1pt,
         labelwidth=0pt,
         listparindent=0pt,
         parsep=0pt}

%------------------------------------------------------------------------------%
% define packages                                                              %
%------------------------------------------------------------------------------%
\usepackage{upref}
\usepackage{amssymb}
\usepackage{slantsc}
\usepackage{amsmath}
\usepackage{amsthm}
\usepackage{thmtools}
\usepackage{amscd}
\usepackage{verbatim}
\usepackage{latexsym}
\usepackage{multicol}
\usepackage{graphicx}
\usepackage{setspace}
\usepackage[ngerman,english]{babel}
\usepackage[protrusion=true,expansion=true]{microtype}
\usepackage{tikz}
\usetikzlibrary{arrows,arrows.meta,decorations.markings}
\usepackage{mathrsfs}
\usepackage{amsfonts}
\usepackage{datetime2}
\usepackage{textgreek}
\usepackage{hyperref}
\usepackage{pifont}
\usepackage{relsize}
%\usepackage{biblatex}
\relsize{-0.05}
\usepackage[skip=0mm]{parskip}
\usetikzlibrary{decorations.pathmorphing}

\usepackage{mlmodern}
\usepackage[T5,T1]{fontenc}
\usepackage{ragged2e}

%\usepackage{mathabx}
%\usepackage{times}

%\usepackage{fontspec}
%\setmainfont{Nimbus Roman No9 L}
%\usepackage[T1]{fontenc}

%------------------------------------------------------------------------------%
% define page layout                                                           %
%------------------------------------------------------------------------------%
\usepackage{fancyhdr}
\pagestyle{fancy}
\setlength{\headheight}{30.0pt}
\renewcommand{\headrulewidth}{0pt}

\AtBeginDocument{
  \addtolength\abovedisplayskip{-0.0\baselineskip}
  \addtolength\belowdisplayskip{-0.0\baselineskip}
}

\fancyhead{}
\fancyfoot{}
\addtolength{\headwidth}{0.02\marginparwidth}

\makeatletter
\let\ps@plain\ps@fancy
\makeatother

\renewcommand\sectionmark[1]
{
 \def\sectionname{#1}
 \markright{\thesection #1}
}

\makeatletter
\@addtoreset{section}{part}
\makeatother

%------------------------------------------------------------------------------%
% define footnote without number                                               %
%------------------------------------------------------------------------------%
\newcommand{\myfootnote}[1]{%
\renewcommand{\thefootnote}{}%
\footnotetext{#1}%
\renewcommand{\thefootnote}{\arabic{footnote}}%
}

%------------------------------------------------------------------------------%
% define numbering in equations                                                %
%------------------------------------------------------------------------------%
\numberwithin{equation}{section}

%------------------------------------------------------------------------------%
% define newtheorem                                                            %
%------------------------------------------------------------------------------%
\declaretheoremstyle[
headfont=\scshape, 
headindent = 0mm, %\parindent,
%postheadhook = {\hspace{0mm}\newline},
%postheadspace = \newline, 
spaceabove = 3mm, 
spacebelow = 3mm,
numberwithin=section,
name=Theorem]{khMATH}
\declaretheorem[style=khMATH]{theorem}

\declaretheoremstyle[
headfont=\bfseries, 
headindent = 0mm, %\parindent,
%postheadhook = {\hspace{0mm}\newline},
%postheadspace = \newline, 
spaceabove = 3mm, 
spacebelow = 3mm,
numberwithin=part,
name=Theorem]{khMATH1}
\declaretheorem[style=khMATH1]{theorem1}

\declaretheoremstyle[
headfont=\scshape, 
headindent = 0mm, %\parindent,
%postheadhook = {\hspace{0mm}\newline},
%postheadspace = \newline, 
spaceabove = 3mm, 
spacebelow = 3mm,
numberwithin=section,
name=Corollary]{khMATH1}
\declaretheorem[style=khMATH1]{corollary}

%            name         counter     label              numeration-mode
\theoremstyle{plain}
\newtheorem*{introduction}            {\sc Introduction}
\newtheorem {standard}                {}                 [section]
\newtheorem*{definition}              {\sc Definition}
\newtheorem{satz1}                    {\sc Satz}
\newtheorem{lemma}                    {\sc Lemma}
\newtheorem*{proposition}             {\sc Proposition}
%\newtheorem{corollary}               {\sc Corollary}
\newtheorem*{corollar}                {\sc Korollar}
\newtheorem*{remark}                  {\sc Remark}
\newtheorem*{remarks}                 {\sc Remarks}
\newtheorem*{example}                 {Example}
\newtheorem*{examples}                {Examples}
\newtheorem*{theo}                    {\sc Theorem}

%------------------------------------------------------------------------------%
% define tabular                                                               %
%------------------------------------------------------------------------------%
\usepackage{tabularx}
\newcolumntype{L}[1]{>{\raggedright\arraybackslash}p{#1}}
\newcolumntype{C}[1]{>{\centering\arraybackslash}  p{#1}} 
\newcolumntype{R}[1]{>{\raggedleft\arraybackslash} p{#1}} 

%------------------------------------------------------------------------------%
% define \sh, \zB                                                              %
%------------------------------------------------------------------------------%
\def\dsh{{d.\,h.}}
\def\ie{{i.\,e.}}
\def\zB{z.\,B.}
\renewcommand{\abstractname}{ }

%------------------------------------------------------------------------------%
% change default spacing for cases                                             %
%------------------------------------------------------------------------------%
\makeatletter
\renewcommand*\env@cases[1][1.6]{%
  \let\@ifnextchar\new@ifnextchar
  \left\lbrace
  \def\arraystretch{#1}%
  \array{@{}l@{\quad}l@{}}%
}
\makeatother

\newenvironment{rcases}
  {\left.\begin{aligned}}
  {\end{aligned}\right\rbrace}

%------------------------------------------------------------------------------%
% change default for \predisplaypenalty                                        %
%------------------------------------------------------------------------------%
\predisplaypenalty=990
\allowdisplaybreaks \numberwithin{equation}{section}

%------------------------------------------------------------------------------%
% general notations                                                            %
%------------------------------------------------------------------------------%
\def\*{\star}
\def\={\,=\,}
\def\+{\,+\,}
\def\xsubset{{\,\subset\,}}
\def\xtimes{{\,\times\,}}
\def\xeof{{\,\in\,}}
\def\eq{{\,=\,}}
\def\xto{{\,\to\,}}
\def\eqdef{\,:=\,}
\def\:{{\,:\,}}
\def\ele{{\,\in\,}}
\def\exists{\mbox{ exists }}
\def\and{\mbox{ and }}
\def\for{\mbox{ for }}
\def\fa{\mbox{ for all }}
\def\fe{\mbox{ for each }}
\def\qfa{\quad \mbox{ for all }}
\def\df{\,\dotfill}
\def\iff{if and only if }
\def\wrt{with respect to}

\makeatletter
\newcommand \Df {\leavevmode \cleaders \hb@xt@ .70em{\hss .\hss }\hfill \kern \z@}
\makeatother

\def\sql{\mathfrak{a}}               % sesquilinear form a
\def\DF{B}                           % Dirichlet form a
%------------------------------------------------------------------------------%
% number sets and special elements                                             %
%------------------------------------------------------------------------------%
\def\P{\mathbb{H}}                   % half plane of $\C$
\def\J{I}                            % interval I
\def\N{{\mathbb{N}}}                 % unsigned integers
\def\Q{{\mathbb{Q}}}                 % rational integers
\def\Z{{\mathbb{Z}}}                 % integers
\def\R{{\mathbb{R}}}                 % real numbers
\def\Rn{{\mathbb{R}}^n}              % real numbers in Rn
\def\Rd{{\mathbb{R}}^d}              % real numbers in Rn
\def\C{{\mathbb{C}}}                 % complex numbers
\def\F{{\mathbb{K}}}                 % arbitrary field
\def\Torus{{\mathbb{T}}}             % complex circle line
\def\T{\tau}                         % upper bound of interval (0,t)
\def\sign#1{\operatorname{sign}(#1)} % sign of a number

%------------------------------------------------------------------------------%
% set theory                                                                   %
%------------------------------------------------------------------------------%
\def\G{G}                            % domain $G$
\def\clG{\overline{\G}}              % closure of $G$
\def\bndG{\partial \G}               % boundary of $G$
\def\setA{\mathcal{A}}               % set A
\def\setB{\mathcal{B}}               % set B
\def\setC{\mathcal{C}}               % set C
\def\setM{M}                         % set M
\def\setN{N}                         % set N
\def\setS{\mathcal{S}}               % set S
\def\famF{\mathfrak{F}}              % family of sets
\def\indI{\mathcal{I}}               % index set I
\def\pot#1{\mathfrak{P}(#1)}         % power set

\def\relR{\mathbf{R}}                % relation R
\def\mapf{f}                         % mapping f
\def\mapg{g}                         % mapping g

\def\set#1#2{\{ \, #1 \,: \, #2 \, \}}
\def\tensor#1#2{#1 \otimes #2}
\def\Tens{\oplus}

\def\cal#1{{\mathcal{#1}}}
\def\aut#1{\operatorname{Aut}(#1)}

\def\cross#1#2{#1 {\,\times\,} #2}
\def\multicross#1#2{#1 \times  \ldots \times  #2}
\def\multitens#1#2 {#1 \otimes \ldots \otimes #2}

%------------------------------------------------------------------------------%
% group theory                                                                 %
%------------------------------------------------------------------------------%
\def\1{{\mathsf{1}}}                % unit element of a monoid
\def\0{{\mathsf{0}}}                % unit element of an Abelian monoid
\def\kernel#1{\mathfrak{N}(#1)}     % kernel of a homomorphism
\def\cokernel#1{\mathfrak{C}{(#1)}} % cokernel of a homomorphism
\def\Hom#1{\operatorname{Hom}(#1)}

%------------------------------------------------------------------------------%
% linear algebra                                                               %
%------------------------------------------------------------------------------%
\def\vecV{\mathit{V}}               % vector space V
\def\vecH{\mathit{H}}               % vector space H
\def\vecW{\mathit{W}}               % vector space W
\def\vecU{\mathit{U}}               % subspace U
\def\convC{\mathcal{C}}             % convex set C
\def\basH{\mathfrak{B}}             % Hamel base B

\def\LO#1{\mathscr{L}(#1) }         % algebra of linear transformations

\def\scalar#1#2{ \langle #1,#2 \rangle }
\def\trace#1{\operatorname{tr}(#1)}
\def\dim#1{{\operatorname{dim} \, #1 }}
\def\codim#1{{\operatorname{codim} \, #1 }}
\def\dim{\operatorname{dim}}
\def\codim{\operatorname{codim}}
\def\ind{\operatorname{ind}}

\def\genG{\cal{G}}

\def\algA{\cal{A}}
\def\algB{\cal{B}}
\def\algU{\cal{U}}

\def\idI{\cal{I}}

%------------------------------------------------------------------------------%
% topology                                                                     %
%------------------------------------------------------------------------------%
\def\compK{{K}}          % compact set C
\def\openO{\mathcal{O}}  % open set O
\def\openU{\mathcal{U}}  % open set U
\def\U{\mathcal{U}}      % neighbourhood U
\def\V{\mathcal{V}}      % neighbourhood V
\def\d{\mathit{d}}       % metric function d

\def\interior#1{\mathring{#1}}
\def\closure#1{{\overline{#1}}}

\def\topT{\mathcal{T}}   % topology T
\def\topX{X}             % topological space X
\def\topY{Y}             % topological space Y
\def\topZ{Z}             % topological space Z
\def\metrS{M}            % metric space M

\def\grpG{G}             % topological group G
\def\grpA{A}             % topological group A
\def\grpB{B}             % topological group B
\def\grpC{C}             % topological group C
\def\grpH{H}             % topological group H
\def\grpU{U}             % topological group U
\def\grpV{V}             % topological group V
\def\topG{G}             % topological group G
\def\topH{H}             % topological group H
\def\topU{U}             % topological subgroup U

\def\subS{\cal{S}}
\def\riR{R}              % ring R

%------------------------------------------------------------------------------%
% calculus                                                                     %
%------------------------------------------------------------------------------%
\def\supp#1{\operatorname{supp}(#1)}
\def\singsupp#1{\operatorname{sing \, supp}(#1)}

\def\re{\operatorname{Re}}
\def\im{\operatorname{Im}}
\def\grad{\operatorname{grad}}

%\def\abs#1 {\left|#1\right|}
\def\abs#1 {\lvert #1 \rvert}
%\def\abs#1 {\left|\,#1\,\right|}

%------------------------------------------------------------------------------%
% measure theory                                                               %
%------------------------------------------------------------------------------%
\def\salgA{\Sigma}               % sigma-algebra S
\def\salgB{\cal{B}}              % sigma-algebra B
\def\Borel#1{\cal{B}_{#1}}       % Borel-algebra 
\def\LB{\lambda}                 % Lebesgue-Borel measure

%------------------------------------------------------------------------------%
% function spaces                                                              %
%------------------------------------------------------------------------------%
\def\Lp{L^p(\G)}                 % spaces L_p
\def\Lq{L^q(\G)}                 % spaces L_q
\def\Lh{L^2(\G)}                 % spaces L_2
\def\Li{L^{\infty}(\G)}          % spaces L_infty
\def\Lc{L_{loc}^1(\G)}           % spaces L_1^{loc}
\def\Ck{C^k(\G)}                 % spaces C_k

\def\BV#1{BV(#1) }               % set of functions of bounded variation
\def\AC#1{AC(#1) }               % set of absolutely continuous functions

\def\MP#1{\mathscr{M}^+(#1) }    % set of positive measures
\def\MS#1{\mathscr{M}^-(#1) }    % vector space of finite signed measures
\def\MC#1{\mathscr{M}(#1) }      % vector space of complex measures
\def\MCR#1{\mathscr{M}_{r}(#1) } % vector space of regular complex measures

\def\SobW{W^{k,p}(\G)}                % Sobolev space W
\def\SobO{\mathring{W}^{k,p}(\G)}     % Sobolev space W
%\def\WN#1#2{\mathring{W}^{#1}(#2) }   % Sobolev space W
\def\WN#1#2{W_0^{#1}(#2) }   % Sobolev space W
%\def\HN#1#2{\mathring{H}^{#1}(#2) }   % Sobolev space H
\def\HN#1#2{H_0^{#1}(#2) }   % Sobolev space H
%------------------------------------------------------------------------------%
% functional analysis                                                          %
%------------------------------------------------------------------------------%
\def\snorm#1{\lVert #1 \rVert}
\newcommand{\norm}[1]{\left\lVert #1 \right\rVert}

\def\topV{V}                       % topological vector space V
\def\topW{W}                       % topological vector space W
\def\normS{X}                      % normed space

\def\banX{\mathit{X}}              % Banach space X
\def\banY{\mathit{Y}}              % Banach space Y
\def\banZ{\mathit{Z}}              % Banach space Z
\def\dbanX{\mathit{X^*}}           % dual space of Banach space X
\def\dbanY{\mathit{Y^*}}           % dual space of Banach space Y
\def\dbanZ{\mathit{Z^*}}           % dual space of Banach space Z

\def\banA{\mathit{A}}              % Banach algebra A
\def\banB{\mathit{B}}              % Banach algebra B
\def\banC{\mathit{C}}              % C*-algebra C


\def\domain#1{\mathfrak{D}(#1)}    % domain
\def\range#1{\mathfrak{R}(#1)}     % range

\def\isoJ{\mathfrak{J}}            % isometry J
\def\repA{{\cal{A}_{\sql}}}        % representation operator A
\def\linS{{S}}                     % linear operator S
\def\linG{{G}}                     % linear operator G
\def\linL{{L}}                     % linear operator L
\def\linT{{T}}                     % linear operator A
\def\linA{\mathit{A}}              % linear operator A
\def\linB{{B}}                     % linear operator B
\def\linM{{M}}                     % linear operator M
\def\linU{{U}}                     % linear operator U
\def\linI{{I}}                     % linear operator I
\def\A{\cal{A}}                    % linear differential operator A
\def\BDO{\cal{B}}                  % linear boundary differential operator B
\def\linU{{U}}                     % unitary operator U
\def\linP{{P}}                     % projection P
\def\compA{k}                      % compact operator k
\def\linFR{f}                      % finite rank operator f
\def\fredA{A}                      % Fredholm operator F

\def\HAM{\cal{H}}                  % Hamilton operator
\def\PA{\partial^{\alpha}}         % elementary differential operator
\def\DO{\cal{A}}                   % linear differential operator
\def\BC{\cal{B}}                   % boundary operator
\def\SLO{\cal{L}_{p,q}}            % Sturm-Liouville operator
\def\MO{\cal{L}_{M}}               % momentum operator

\def\HS#1{S(#1)           }        % algebra of Hilbert-Schmidt operators
\def\FR#1{F(#1)           }        % algebra of finite rank operators
\def\TC#1{N(#1)           }        % algebra of trace class operators
\def\BU#1{\mathscr{U}(#1) }        % algebra of unitary linear operators
\def\BO#1{\mathscr{L}(#1) }        % algebra of bounded linear operators
\def\CO#1{\mathscr{C}(#1) }        % algebra of compact linear operators
\def\CL#1{\mathscr{C}(#1) }        % algebra of compact linear operators
\def\FO#1{\Phi(#1) }               % set of Fredholm operators
\def\FK#1#2{\Phi_{#2}(#1) }        % set of Fredholm operators with index k

\def\hilH{\mathit{H}}              % Hilbert space H
\def\clU{\mathit{U}}               % closed subspace U

\def\orthO{\mathfrak{O}}           % orthonormal system O
\def\orthS{\mathfrak{S}}           % orthonormal system S
\def\basB{\mathfrak{B}}            % basis B of a Hilbert space

\def\GL#1 {\mathbf{GL}(#1)}        % linear group
\def\OG#1 {\mathbf{O}(#1)}         % linear group
\def\SOG#1 {\mathbf{SO}(#1)}       % linear group

%------------------------------------------------------------------------------%
% distribution theory                                                          %
%------------------------------------------------------------------------------%
\def\disF{F}                  % distribution F
\def\disG{G}                  % distribution G
\def\disH{H}                  % distribution H
\def\disU{U}                  % distribution U
\def\disT{T}                  % distribution T
\def\SF{C^{\infty}(\G)}       % space of smooth functions
\def\TF{C_c^{\infty}(\G)}     % space of compactly supported smooth functions
\def\TFRd{C_c^{\infty}(\R^d)} % space of test functions
\def\SRd{\mathscr{S}(\R^d)}   % Schwartz space
\def\SR1{\mathscr{S}(\R)}     % Schwartz space
\def\B1#1{C^1(#1) }           % space of continuously differentiable functions
\def\Bm#1{C^m(#1) }           % space of continuously differentiable functions
\def\Ck{C^k(\G)}
\def\TD{\mathscr{S}'(\R^d)}   % space of temperated distributions
\def\E{\cal{E}}               % space of distributions with compact support
\def\D{\mathscr{D}}           % D
\def\FT{\mathcal{F}}          % Fourier transformation F
\def\FT{\mathscr{F}}          % Fourier transformation F

%------------------------------------------------------------------------------%
% other definitions                                                            %
%------------------------------------------------------------------------------%
\def\Proof{\textsc{Proof}}
\def\FRECHET{Fr\'{e}chet}
\def\ARZELA{Arzel\'{a}}
\def\GARDING{G\r{a}rding}
\def\HOELDER{H\"older}
\def\SCHROEDINGER{Schr\"odinger}
\def\JOERGENS{J\"orgens}
\def\WUEST{W\"uest}
\def\KAEHLER{K\"aehler}
\def\HOEFLINGER{H\"oflinger}
\def\POINCARE{Poincar\'{e}}
\def\FA{Functional Analysis}

\def\Author   {K.\,{\HOEFLINGER}}
\def\AuthorUC {K.\,HOEFLINGER}
\def\Version  {1.0}
\def\Email    {khoeflinger@t-online.de}


%\renewcommand\thesection{\Roman{section}}

\def\Title {Fredholm Operators and the Essential Spectrum}

\titleformat{\section}
{\normalfont \bfseries \filcenter}
{\thesection.\;}{0mm}{}

\titleformat{\subsection}%[runin]
{\normalfont \itshape}
{}{}{}

\titlespacing*{\section}   {0pt}{12pt}{4pt}
\titlespacing*{\subsection}{0pt}{6pt}{0pt}

%------------------------------------------------------------------------------%
%                                BEGIN OF DOCUMENT                             %
%------------------------------------------------------------------------------%
\begin{document}
\thispagestyle{empty}

\centerline{\large{\textbf{\Title}}}

\vspace{4pt}
\centerline{\footnotesize{{\textsc{Martin Schechter}}}}
\vspace{6pt}

\fancyhead[C] {\small{\textsc{\Title}}}
\fancyhead[OL]{\small{\thepage}}
\fancyhead[ER]{\small{\thepage}}

\myfootnote
{
  Prepared under contract no. DA31-124-ARO- D-204
  with the U.S. Army research office.
  January 1965
}

\begin{abstract}
Fredholm operators are defined and their basic properties are proved.
The essential spectrum of an arbitrary closed operator is considered
and its invariance under several kinds of perturbation is proved.
\end{abstract}

%------------------------------------------------------------------------------%
\section{Introduction}                                                         %
%------------------------------------------------------------------------------%
We first give a complete treatment of Fredholm operators.
Most of the results are found in the literature (e.g., [4,8]),
but we have been able to substantially simplify proofs in several instances.
The deepest theorem we make use of is the closed graph theorem.
We are able to avoid completely the concept of the opening between two subspaces.

\subparagraph*{}
Next we give a proof of a classical theorem of Weyl [13]
on the invariance of the essential spectrum under compact perturbations
and give two definitions for general closed operators in a Banach space,
one of them due to Wolf [14].
We then consider more general types of perturbations.
Further results may be found in [10].

\subparagraph*{}
Remarks on specific results and methods are given at the end.

%------------------------------------------------------------------------------%
\section{Fredholm Operators}                                                   %
%------------------------------------------------------------------------------%
A linear operator $A$ from a Banach space $X$ to a Banach space $Y$
will be called a Fredholm operator if

\begin{itemize}
\item[1.] $A$ is closed;
\item[2.] the domain $\domain{A}$ of $A$ is dense in $X$;
\item[3.] $\alpha(A)$, the dimension of the null space $\kernel{A}$ of $A$,
          is finite;
\item[4.] $\range{A}$, the range of $A$, is closed in $Y$;
\item[5.] $\beta(A)$, the codimension of $\range{A}$ in $Y$, is finite.
\end{itemize}

\paragraph*{}
The index of a Fredholm operator $A$ is defined as
\[
i(A) = \alpha(A) - \beta(A).
\]
We now make several observations.
The set of Fredholm operators from $X$ to $Y$ will be denoted by $\Psi(X,Y)$.

\paragraph*{}
\begin{itemize}[leftmargin=12mm]
\item[(O-1)]
If we decompose $X$ into
\begin{equation}\tag{1}
X = N(A) \oplus X',
\end{equation}
where $X'$ is a closed subspace of $X$, then
\begin{equation}\tag{2}
\norm{x} \, \leq \, \mbox{const.} \norm{A \,x}
\end{equation}
holds for all $x \in \domain{A} \cap X'$. 
This merely expresses the fact
that $A$ restricted to $\domain{A} \cap X'$
has an inverse defined everywhere on $\range{A}$,
which is a Banach space.
Hence this closed inverse is bounded.

\item[(O-2)]
If $A$ is closed, $\alpha(A) < \infty$ and (2) holds,
then $\range{A}$ is closed.
For if $A x_n \to y$,
$x_n \in \domain{A} \cap X'$,
then $x_n \to x \in X'$.
Thus $x \in \domain{A}$ and $Ax = y$.

\item[(O-3)]
If $A \in \Psi(X,Y)$,
there is a bounded operator $A'$ from $Y$ to $\domain{A} \cap X'$ such that
\begin{equation}\tag{3}
A'A = I \mbox{ on } \domain{A} \cap X'
\end{equation}
\begin{equation}\tag{4}
AA' = I \mbox{ on } \range{A},
\end{equation}
while $A'$ vanishes on any prescribed complement of $\range{A}$ in $Y$.
(Here $I$ denotes the identity operator.)

\paragraph*{}
\Proof.
Obvious.

\item[(O-4)]
The operator $A'$ also satisfies
\begin{equation}\tag{5}
A'A = I +F_1 \mbox{ on } \domain{A}
\end{equation}
\begin{equation}\tag{6}
AA ' = I + F_2 \mbox{ on } Y
\end{equation}
where the operators $F_i$ are bounded and have finite dimensional
ranges.
In fact $A'A -I$ vanishes on $\domain{A} \cap X'$
and hence maps $\domain{A}$ into the image of $N(A)$.
A similar statement holds for $AA' -I$.

\item[(O-5)]
If $A$ is closed, $\domain{A}$ dense and there are bounded
operators $A_i$ from $Y$ to $X$ and compact operators $K_i$ on $X$,
$K_2$ on $Y$
such that
\begin{equation}\tag{7}
A_1 A = I + K_1 \mbox{ on } \domain{A}
\end{equation}
\begin{equation}\tag{8}
A A_2 = I + K_2 \mbox{ on } Y,
\end{equation}
then $A \in \Psi(X,Y)$.

\vspace{3mm}

\Proof.
$\alpha(A) \leq \alpha(I + K_1) < \infty$ and
$\beta(A)  \leq \beta(I + K_2)  < \infty$
(F. Riesz).
We must show that $\range{A}$ is closed.
Suppose $\{x_n\} \subset \domain{A} \cap X'$,
$\norm{x_n} \eq 1$ and $A x_n \to 0$.
Since $A$ is bounded,
$(I + K_1)x_n \to 0$.
Since $\{x_n\}$ is bounded,
there is a subsequence (also denoted by $\{x_n\}$)
for which $K_1 x_n$ converges to some element $z$.
Thus $x_n \to -z$.
Since $A$ is closed, $z \in \domain{A}$ and $Az = 0$,
i.e. $z \in \kernel{A}$.
Since $X'$ is closed, $x \in X'$ and hence $z = 0$.
But $\norm{z} = \lim \norm{x_n} \eq 1$.
This shows that (2) holds.
Thus $\range{A}$ is closed by (O-2).
\end{itemize}

\paragraph*{}
We shall also need the following important lemma.

\begin{lemma} % lemma 1
If $A \in \Psi(X,Y)$ and $B \in \Psi(Z,X)$, then $AB \in \Psi(Z,Y)$ and
\begin{equation}\tag{9}
i(AB) = i(A) + i(B) .
\end{equation}
\end{lemma}

Proof.
Clearly, $\domain{AB}$ is dense in $Z$ (O-1).
Set
\[
X_0 = \range{B} \cap \kernel{A}
\]
\[
\range{B} = X_0 \oplus X_1
\]
\[
\kernel{A} = X_0 \oplus X_2
\]
\[
X = \range{B} \oplus X_2 \oplus X_3.
\]
Then $X_0,X_2,X_3$ are finite dimensional and $X_1$ is closed.
Since $\domain{A}$ is dense in $X$,
we can take $X_3$ to be contained in $\domain{A}$.
Set $d_i = \dim X_i$.
Then
\[
\alpha(AB) = \alpha(B) + d_0
\]
\[
\beta(AB) = \beta(A) + d_3
\]
\[
d_0 + d_2 = \alpha(A)
\]
\[
d_2 + d_3 = \beta(B).
\]
These show that $\alpha(AB)$ and $\beta(AB)$ are finite and (9) holds.
We must also show that $\range{AB}$ is closed.
Now $\range{AB}$ is the range of $A$ on $\domain{A} \cap X_1$.
If $A x_n \to f$ for $\{x_n\} \subset \domain{A} \cap X_1$,
we have by (2)
\[
\norm{x_m-x_n} \leq \mbox{ const. } \norm{A(x_m-x_n} \to 0
\]
and since $X_1$ is closed, $x_n \to x \in X_1$.
Since $A$ is closed,
$x \in \domain{A}$ and $Ax = f$.
Thus $\range{AB}$ is closed.
It remains to show that $AB$ is closed.
Suppose $\{z_n\} \subset \domain{AB}$, $z_n \to z$, $AB z_n \to y$.
Write
\[
Bz_n = x_n^{(0)} + x_n^{(1)}
\]
where $x_n^{(0)} \in \kernel{A}$ and $x_n^{(1)} \in X'$.
Then by (2) $x_n^{(1)} \to x^{1} \in X'$.
If $\snorm{x_n^{(0)}} \leq \mbox{ const.}$,
there is a subsequence
(also denoted by $\{x_n\}$)
for which $x_n^{(0)} \to x^{(0)} \in \kernel{A}$.
Thus $Bz_n \to x^{(0)} + x^{(1)}$,
and since $B$ is closed, $z \in \domain{B}$ and $Bz = x^{(0)}+x^{(1)}$.
Since $A$ is closed $x^{(0)} + x^{(1)} \in \domain{A}$
and $A(x^{(0)} + x^{(1)}) = y$.
Thus $z \in \domain{AB}$ and $ABz=y$.
We now show that $\{x_n^{(0)}\}$ is bounded.
For if $\zeta_n = \snorm{x_n^{(0)}} \to \infty$,
set $u_n = \zeta_n^{-1} x_n^{(0)}$.
Then $\norm{u_n} \eq 1$ and hence there is a subsequence
(also denoted by $\{u_n\}$)
such that $u_n \to u \in \kernel{A}$.
Moreover
\[
B(\zeta_n^{-1} z_n) -u_n = 
\zeta_n^{-1} (Bz_n - x_n^{(0)}) = 
\zeta_n^{-1} x_n^{(1)}) \to 0. 
\]
Hence $B(\zeta_n^{-1} z_n) \to u$.
Since $\zeta_n^{-1} z_n \to 0$ and $B$ is closed,
$u \eq 0$.
This is impossible since $\norm{u} = \lim \norm{u_n} \eq 1$.
This completes the proof.

\begin{theorem1} % theorem 1
If $A \in \Psi(X,Y)$ and $K$ is a compact operator from $X$ to $Y$,
then $A+K \in \Psi(X,Y)$ and $i(A+K) = i(A)$.
\end{theorem1}

Proof.
By (O-4) there is a bounded operator $A'$ such that (5) and (6) hold.
Now
\[
A'(A + K) = I+F_1 +A'K \mbox{ on } \domain{A}
\]
\begin{equation}\tag{10}
(A+K)A' = I +F_2 +KA' \mbox{ on } Y.
\end{equation}
Since $A'$ is bounded, the operators $A'K$ and $KA'$ are compact.
We now apply (O-5) to obtain that $A + X \in \Psi(X,Y)$.
Now $A'$ is a bounded operator from $Y$ to $\domain{A}$.
If we equip $\domain{A}$ with the graph norm of $A$,
it becomes a Banach space $\tilde{X}$
and $A'$ becomes a Fredholm operator from $Y$ to $\tilde{X}$.
Also $A \in \Psi(\tilde{X},Y)$ with $\alpha(A)$ and $\beta(A)$ the same.
Hence by (5)
\begin{equation}\tag{11}
i(A) + i(A') = i(I+F_1) = 0 \quad \mbox{(Riesz)}.
\end{equation}
But by (10)
\[
i(A+K) + i(A') = i(I+F_2 + KA') = 0.
\]
Hence $i(A+K) = i(A)$ and the proof is complete.

\begin{theorem1} % theorem 2
For $A \in \Psi(X,Y)$ there is an $\epsilon > 0$
such that
for any bounded operator $T$ from $X$ to $Y$
with $\norm{T} < \epsilon$ one has
$A + T \in \Psi(X,Y)$, $i(A+T) \eq i(A)$ and $\alpha(A+T) < \alpha(A)$.
\end{theorem1}

Proof.
By (O-4) there is a bounded operator $A'$ satisfying (5) and (6).
Then
\[
A'(A + T) = I +F_1 +A'T
\]
\begin{equation}\tag{12}
(A+T)A' = I+F_2+TA'.
\end{equation}
Take $\epsilon = \norm{A'}^{-1}$.
Then $\norm{A'T} < 1$ and $\norm{TA'} < 1$.
Thus the operators $I +A'T$ and $I+TA'$ are invertible and
\[
(I + A'T)^{-1}A'(A+T) = I + (U+A'T)^{-1}F_1
\]
\[
(A + T)A'(1+TA')^{-1} = I + F_2(I+TA')^{-1}
\]
Thus by (O-5), $A+T \in \Psi$.
By (12) and theorem 1
\[
i(A+T) + i(A') = i(I+F_2+TA') = i(1+TA') = 0.
\]
Combining this with (11) we see that $i(A+T) = i(A)$.
It remains to prove $\alpha(A + T) \leq \alpha(A)$.
By (O-3),
$A'(A + T) = I + A'T$ on $\domain{A} \cap X'$
and hence is one-to-one on this manifold.
Moreover $\kernel{A+T} \cap X' = \{0\}$.
For if $y$ is in this set,
it is in $\domain{A} \cap X'$ and $(A+T)y = 0$.
Thus $(I+A'T)y = 0$ showing that $y = 0$.
Since 
\[
\kernel{A+T} \oplus X' \subseteq X \= \kernel{A} \oplus X',
\]
we see that $\dim \kernel{A +T} \leq \dim \kernel{A}$
and the proof is complete.

\paragraph*{}
In the sequel we shall make continual use of the following lemma:

\begin{lemma} % lemma 2
Let $\tilde{X}$ be a Banach space dense in $X$
and let $A$ be a linear operator from $X$ to $Y$
with $\domain{A} = \tilde{X}$.
If $A \in \Psi(\tilde{X},Y)$, then $A \in \Psi(X,Y)$.
\end{lemma}

Proof.
Let $E$ be the linear operator from $X$ to $\tilde{X}$ with
$\domain{E} = \domain{A}$ and $Ex = x$ for $x \in \domain{E}$.
One checks easily that $E \in \Psi(X,\tilde{X})$.
Thus $AE \in \Psi(X,Y)$ by lemma 1.
But $A = AE$ and the conclusion follows.

\paragraph*{}
An operator $B$ from $X$ to $Y$ is called $A$-compact
if
$\domain{B} \supset \domain{A}$ and
for every sequence $\{x_n\} \subset \domain{A}$ such that
\[
\norm{x_n} + \norm{Ax_n} \leq \mbox{const.}
\]
the sequence $\{Bx_n\}$ has a convergent subsequence.

\begin{theorem1} % theorem 3
If $A \in \Psi(X,Y)$ and $B$ is $A$-compact,
then $A+B \in \Psi(X,Y)$ and $i(A+B) = i(A)$.
\end{theorem1}

Proof.
If we equip $\domain{A}$ with the graph norm $\norm{x}+\norm{Ax}$,
it becomes a Banach space $\tilde{X}$ (since $A$ is closed).
If we now consider $A$ as an operator from $\tilde{X}$ to $X$,
we have $A \in \Psi(\tilde{X},Y)$.
Moreover $B$ is compact from $\tilde{X}$ to $Y$.
Hence by theorem 1, $(A+B) \in \Psi(\tilde{X},Y)$ with
$i(A+b) \eq i(A)$.
We now apply lemma 2 to obtain $(A+B) \in \Psi(X,Y)$.

\begin{theorem1} % theorem 4
For each $A \in \Psi(X,Y)$ there is an $\epsilon > 0$
such that
\[
\norm{B x} \leq \epsilon (\norm{x} + \norm{A x}), \quad x \in \domain{A},
\]
holding for any operator $B$ from $X$ to $Y$
with $\domain{B} \supset \domain{A}$
implies that
$A+B \in \Psi((X,Y)$, $i(A+B) = i(A)$,
and $\alpha(A + B) \leq \alpha(A)$.
\end{theorem1}

Proof.
Similar to that of theorem 3.

\paragraph*{}
For an arbitrary operator $A$ acting in a Banach space $X$
the Fredholm set of $A$,
denoted by $\Phi_A$,
is the set of those complex $\lambda$
for which $A-\lambda \in \Psi(X,X)$.
We have by theorem 1 and 2

\begin{theorem1} % theorem 5
The set $\Phi_A$ is open and $i(A-\lambda)$ is constant
on each component.
\end{theorem1}

\begin{theorem1} % theorem 6
If $K$ is a compact operator in $X$,
then $i(A+K-\lambda) = i(A-\lambda)$ there.
\end{theorem1}

\paragraph*{}
In addition one also has

\begin{theorem1} % theorem 7
On each component of $\Phi_A$,
the quantities
$\alpha(A - \lambda)$ and $\beta(A - \lambda)$ 
are constant except possibly on a discrete set of points
where they take on larger values.
\end{theorem1}

\paragraph*{}
The proof of theorem 7 can be made to rest upon lemma 3.

\begin{lemma} % lemma 3
If $A$ and $B$ are linear operators from $X$ to $Y$
with $A \in \Psi(X,Y)$ and $B$ bounded,
then there is an $\epsilon > 0$
such that
$\alpha(A-\lambda B)$ is constant for $0 < \abs{\lambda} < \epsilon$.
\end{lemma}

Proof.
Assume first that $i(A) = 0$.
Let $x_1,\ldots,x_k$ be a basis for $\kernel{(A}$ and
$y_1,\ldots,y_k$ be a basis for some direct complement of $\range{A}$ in $Y$.
Let $x'_1,\ldots,x'_k$ be a system of bounded linear functionals on $X$
such that
\[
x'_i(x_j) = \delta_{ij} \quad \mbox{(= Kronecker delta)}.
\]
One checks easily that the operator defined by
\[
\tilde{A}x = Ax + \sum_{j=1}^k x'_j(x) y_j
\]
is continuously invertible.
Now $x$ is a solution of $(A-\lambda B)x =0$
if and only if
\[
(\tilde{A} - \lambda B)x = \sum_{i=1}^k x'_i(x) y_i.
\]
For some $\epsilon_1 > 0$, $\tilde{A}- \lambda B$ is continuously invertible
for $\abs{\lambda} < \epsilon_1$ and hence
\[
x = \tilde{A}^{-1} \sum_{j=1}^{\infty} \lambda^j(B \tilde{A}^{-1})^j
\sum_{i=1}^k x'_i(x) y_i
\]
\[
= \sum_{i=1}^k x'_i(x)
\cdot \sum_{j=0}^{\infty} \lambda^j (\tilde{A}^{-1}B)^j \, x_i
\]
\[
= \sum_{i=1}^k x'_i(x) f_i(\lambda),
\]
where $f_i(\lambda)$ are known vector valued functions in $X$ analytic
on $\abs{\lambda} < \epsilon_1$.
Thus
\[
x'_m(x) = \sum_{i=1}^k \, x'_i(x) x'_m(f_i(\lambda))
\]
\[
\sum_{i=1}^k \, [\delta_{mi}-x'_i(f_i(\lambda)) ] \, x'_i(x) = 0.
\]
Conversely, if $\xi_1,\ldots,\xi_k$ are solutions of
\begin{equation}\tag{15}
\sum_{i=1}^k [\delta_{mi}-x'_n(f_i(\lambda)) ] \xi = 0,
\end{equation}
then by working back we see that
\[
x \= \sum_{i=1}^k \xi_i \, f_i(\lambda)
\]
is a solution of $(A-AB)x = 0$.
Thus $\alpha(A-AB)$ coincides with
the number of linearly independent solutions of (13).
If every coefficient in (13) vanishes identically
in $\abs{\lambda} < \epsilon_1$,
there are exactly $k$ linearly independent solutions and
$\alpha(A-AB) = k$ for $\abs{\lambda} < \epsilon_1$.
Otherwise there is a minor of largest order in the determinant of (13)
which does not vanish identically.
Since this minor is analytic in $\lambda$,
it can vanish at most at isolated points.
Thus there is an $\epsilon > 0$
such that this minor is different from $0$ in $0 < \abs{\lambda} < \epsilon$.
The number of independent solutions of (13) is constant on this set
and hence the same is true for $\alpha(A-AB)$.
Thus the lemma is proved for $i(A) = 0$.

\subparagraph*{}
Next consider the case $i(A) > 0$.
Let $Z$ be a normed $i(A)$-dimensional space and let $Y_1 =Y + Z$ with norm
$\abs{y + z} = \abs{y} \oplus \abs{z} $ when $y \in Y$, $z \in Z$.
Consider $A$ and $B$ operators from $X$ to $Y_1$.
For these spaces $i(A) = 0$ and the previous case applies.
But $\alpha(A- \lambda B)$ is not changed when we replace $Y$ by $Y_1$.
Thus the lemma is proved in this case.

\subparagraph*{}
If $i(A) < 0$, let $Z$ be an $i(A)$-dimensional space
and set $X_1 = X \oplus Z$.
Let $\hat{A}$ and $\hat{B}$ be extensions of $A$ and $B$ to $X$,
which vanish on $Z$.
Then $i(\hat{A}) = 0$ and the first case applies.
We then merely observe that
$\alpha(\hat{A} - \lambda \hat{B}) = \alpha(A-\lambda B) + i(A)$
and the proof is complete.

\paragraph*{}
Proof of theorem 7.
Let $\Omega$ be a component of $\Phi_A$.
Let $\lambda_0$ be any point in $\Omega$
where $\alpha(A-\lambda)$ has its minimum value in $\Omega$.
Then by lemma 3 there is a neighbourhood about $\lambda_0$
where $\alpha(A-\lambda)$ has this constant value.
Let $\lambda_1$ be any other point in $\Omega$.
It suffices to prove that there is a deleted neighbourhood about
$\lambda_1$ in which $\alpha(A-\lambda)$ has this minimum value.
Connect $\lambda_1$ to $\lambda_0$ by a smooth curve in $\Omega$.
At each point of this curve there is a deleted neighbourhood
in which $\alpha(A-\lambda)$ is constant (Lemma 3).
Since the curve is compact,
there is a finite set of such neighbourhoods covering it.
Since in each deleted neighbourhood $\alpha(A-\lambda)$ is constant
and each one overlaps with at least one other,
they all have the same constant value for $\alpha(A-\lambda)$,
namely the minimum value.
Q.E.D.

\paragraph*{}
In later applications we shall need

\begin{lemma} % lemma 4
If $B \in \Psi(Z,X)$, $A$ is an operator from $X$ to $Y$
with $\domain{A}$ dense in $X$,
and $AB \in \Psi(Z,Y)$, then $A \in \Psi(X,Y)$.
\end{lemma}

Proof.
Let $B'$ be the operator corresponding to $B$ described in (O-3).
One easily checks that $B' \in \Psi(X,\tilde{Z})$
where $\tilde{Z}$ is $\domain{B}$ under the graph norm of $B$.
By hypothesis $AB \in \Psi(\tilde{Z},Y)$ and hence $ABB' \in \Psi(X,Y)$.
But $ABB' \eq A+AF$,
where $F$ is bounded and of finite rank (O-4).
Hence $AF$ is a compact operator.
We now apply theorem 1 to conclude that $A \in \Psi(X,Y)$.

\subparagraph*{}
\begin{lemma} % lemma 5
If $A \in \Psi(X,Y)$, $B$ is an operator from $Z$ to $X$
with $\domain{B}$ dense in $Z$, and $AB \in \Psi(Z,Y)$,
then $B \in \Psi(Z,X)$.
\end{lemma}

Proof.
Following the same procedure as in the proof of lemma 4,
we have $A'AB \in \Psi(Z,X)$, and hence in $\Psi(\tilde{Z},X)$.
Since $A'AB = A+FB$,
where $F$ is of finite rank,
the operator $FB$ is compact from $\tilde{Z}$ to $X$.
Hence $B \in \Psi(\tilde{Z},X)$ (theorem 1)
and we can apply lemma 2 to conclude that $B \in \Psi(Z,X)$.

%------------------------------------------------------------------------------%
\section{Weyl's Theorem and Generalizations}                                   %
%------------------------------------------------------------------------------%
Let $T$ be a self-adjoint operator in a Hilbert space $H$
and let $K$ be a symmetric, completely continuous operator in $H$.
Then $T+K$ is also self-adjoint.
One might ask how $\sigma(T)$ and $\sigma(T+K)$ compare.
The answer was given by H. Weyl [13]:

\begin{theorem1} % theorem 8
If $\lambda$ is in $\sigma(T)$ but not in $\sigma(T+K)$,
it must be an isolated eigenvalue of finite multiplicity.
\end{theorem1}

Proof.
We note first that since $Z$ is self-adjoint,
it is closed, and hence $\rho(T) \subseteq \Phi_T$.
Thus if $\lambda \in \sigma(T)$ but not in $\sigma(T+K)$,
it must be in $\Phi_T = \Phi_{T+K}$,
for otherwise it could not be in $\rho(T+K)$.
Thus $T - \mu$ is a Fredholm operator for $\mu$ in a neighbourhood of $\lambda$.
Thus $\alpha(T-\mu)$ is constant in some deleted neighbourhood of $\lambda$
(lemma 3).
Since all non-real $\mu$ are in $\rho(T)$,
we have $\alpha(T-\mu) \eq 0$ in this deleted neighbourhood.

\begin{theorem1} % theorem 9
If $\lambda \in \sigma(T)$ is an isolated eigenvalue
of finite multiplicity,
there is a symmetric, completely continuous operator $K$ in $H$
such that $\lambda \in \rho(T+K)$.
\end{theorem1}

Proof.
We first show that $\lambda \in \Phi_T$.
Since $\beta(T-A) = \alpha(T- A) < \infty$,
the only thing which must be verified
is that $\range{T-A}$ is closed in $H$,
i.e.,
that the inequality
\begin{equation}\tag{14}
\norm{x} \leq C \, \norm{(T-A)x}
\end{equation}
holds for $x \in \domain{T} \cap \kernel{T-A}^{\bot}$ (O-1).
Now $\lambda$ is the only point of $\sigma(T)$ in some interval $[a,b]$,
where $a < \lambda < b$.
If $\{E_{\lambda}\}$ is the spectral family for $T$,
then the projection $E_b - E_a$ maps into $\kernel{T-A}$.
But
\[
\norm{(T-\lambda)x}
= \int_{-\infty}^{\infty} (\mu-\lambda)^2 \, d\norm{E_{\mu}x}^2
\]
\[
\geq (b-\lambda)^2 \int_b^{\infty} d\norm{E_{\mu}x}^2
+    (a-\lambda)^2 \int_{-\infty}  d\norm{E_{\mu}x}^2
\]
\[
\= (b-\lambda)^2 \norm{(1-E_b)x}^2 + (a-\lambda)^2 \norm{E_a x}^2.
\]
If $x \in \domain{T} \cap \kernel{T-A}^{\bot}$,
$(E_b - E_a)x = 0$ and hence (14) holds.
Now let $h_1,\ldots,h_k$ be a basis for $\kernel{T-A}$ satisfying
\[
\scalar{h_i}{h_j} = \delta_{ij}
\]
and set
\[
Kx = \sum_{j=1}^k \scalar{x}{h_j} \, h_j.
\]
Then $K$ is bounded, symmetric and of finite rank.
Thus $T+K-\lambda$ is a Fredholm operator.
Since $\alpha(T+K-\lambda) = \beta(T+K-\lambda) = 0$,
$\lambda \in \rho(T+K)$.

\paragraph*{}
We see from Weyl's theorem that the points of $\sigma(T)$
which remain invariant under any symmetric, completely continuous
perturbations are precisely those points which are not isolated
eigenvalues of finite multiplicity.
The set of such points is called the essential spectrum of $T$
and denoted by $\sigma_e(T)$.
Thus theorems 8 and 9 can be written as

\begin{theorem1} % theorem 10
If $T$ is self-adjoint and $K$ is symmetric and completely continuous,
then $\sigma_e(T + K) = \sigma_e(T)$.
\end{theorem1}

\subparagraph*{}
From the proofs of theorems 8 and 9 we also have

\begin{corollary} % corollary 1
$\sigma_e(T)$ is the complement of $\Phi_T$ in the complex plane.
\end{corollary}

For many applications it is important to generalize to
arbitrary closed operators in Banach space and arbitrary
compact perturbations.
It is no longer true that the invariant points of the spectrum are those
which are not isolated eigenvalues of finite multiplicity.
Wolf [14] defines the essential spectrum
by means of property expressed in Corollary 1.
We denote the essential spectrum according to this definition by
$\sigma_{ew}(A)$.

\begin{definition}
For any closed operator $A$ in a Banach space
$\sigma_{ew}(A) = \C \setminus \Phi_A$.
\end{definition}

\subparagraph*{}
We have immediately by theorem 6:

\begin{theorem1} % theorem 11
For any completely continuous operator $K$,
\[
\sigma_{ew}(A+K) \= \sigma(A).
\]
\end{theorem1}

We shall also employ another definition of essential spectrum:

\begin{definition}
For a closed operator $A$ in a Banach space $X$,
$\sigma_{em}(A)$ is the largest subset of $\sigma(A)$
which remains invariant under arbitrary compact perturbations.
\end{definition}

\begin{theorem1} % theorem 12
$\sigma_{em}(A)$ consists of all complex $\lambda$
except of those $\lambda \in \Phi_A$ with $i(A- \lambda) \eq 0$.
\end{theorem1}

Proof.
If $\lambda \in \Phi_A$ and $i(A-\lambda) \eq 0$,
then the example given in the proof of lemma 3 shows
that there is an operator $F$ of
finite rank for which $\lambda \in \rho(A + F)$.
If $i(A - \lambda) \neq 0$,
then for every compact operator $K$, $i(A+K-\lambda) \neq 0$ (theorem 6)
and hence $\lambda \in \sigma(A + K)$.
If $\lambda \notin \Phi_A$,
then $\lambda \notin \Phi_{A+K}$ and again $\lambda \in \sigma(A+K)$.

\paragraph*{}
We shall now give some theorems on the invariance of
$\sigma_{ew}(A)$ and $\sigma_{em}(A)$ under different types of perturbations,
Throughout we shall assume that $A$ is a closed operator
with dense domain in a Banach space $X$
and $B$ an operator in $X$ with $\domain{B} \supset \domain{A}$.

\begin{theorem1} % theorem 13
If $B$ is $A$-compact, then
\begin{equation}\tag{15}
\sigma_{ew}(A+B) \= \sigma_{ew}(A)
\end{equation}
and
\begin{equation}\tag{16}
\sigma_{em}(A+B) \= \sigma_{em}(A)
\end{equation}
\end{theorem1}

Proof.
We prove the theorem by showing that $\Phi_A = \Phi_{A+B}$
and that $i(A-\lambda) = i(A+B-\lambda)$ for $\lambda \in \Phi_A$.
If $\lambda \in \Phi_A$, then
$\lambda \in \Phi_{A+B}$ and the index relationship holds (theorem 3).
Hence $\Phi_A \subseteq \Phi_{A+B}$.
If $\lambda \in \Phi_{A+B}$, it follows, in particular,
that $A+B$ is closed.
Since $A$ is closed, we have
\begin{equation}\tag{17}
\norm{Ax} \leq C(\norm{x} + \norm{(A+B)x}),
\end{equation}
which expresses the fact that $A$ is a closed operator defined everywhere
on the Banach space $\domain{A+B} \eq \domain{A}$ 
under the graph norm $\norm{x} + \norm{(A+B)x}$.
Hence it is bounded on this set.
From (17) it follows that $B$ is $(A+B)$-compact.
We now apply theorem 3 to $A+B$ with perturbation $-B$.
Thus $\lambda \in \Phi_A$.
Hence $\Phi_{A+B} = \Phi_A$,
and the theorem is proved.

\begin{theorem1} % theorem 14
If $\lambda \in \rho(A) \, \cap \, \Phi_{A+B}$ and either
$(A-\lambda)^{-1}B$ or $(A - \lambda)^{-1}$ is $A$-compact then (15) holds.
If, In addition,
$i(A+B-\lambda) = 0$, then (16) holds.
\end{theorem1}

Proof.
We employ the identities
\begin{equation}\tag{18}
(A+B-\mu) - (A-\mu)(A- \lambda)^{-1}(A+B-\lambda) =
(\mu -\lambda)(A-\lambda)^{-1}B
\end{equation}
\begin{equation}\tag{19}
(A+B-\mu) - (A+B - \lambda)(A- \mu)(A- \lambda)^{-1} =
(\mu - \lambda)B(A - \lambda)^{-1}.
\end{equation}
Assume that $(A-\lambda)^{-1}B$ is $A$-compact.
We equip $\domain{A}$ with the norm $\norm{x} + \norm{Ax} + \norm{Bx}$.
Under it,
$\domain{A}$ becomes a Banach space $\tilde{X}$.
For if $x_n \to X$, $Ax_n \to y$, $Bx_n \to z$,
then $x \in \domain{A}$ and $Ax = y$ since $A$ is closed,
while $(A+B)x = y+z$ because $A+B$ is closed.
Therefore $Bx = z$.
Now if $\mu \in \Phi_{A+B}$, $A+B - \mu$ is a Fredholm operator
from $\tilde{X}$ to $X$.
Moreover,
$(A - \lambda)B$ is a compact operator from $\tilde{X}$ to $X$.
Thus
\begin{equation}\tag{20}
(A- \mu)(A- \lambda)^{-1}(A+B-\lambda)
\end{equation}
is a Fredholm operator from $\tilde{X}$ to $X$ (theorem 1).
Now
$(A-\lambda)^{-1}(A+B-\lambda) \in \Psi(\tilde{X},\tilde{X})$ (Lemma 1).
Thus $A-\mu \in \Psi(\tilde{X},X)$ (Lemma 4)
and hence $A - \mu \in \Psi(X,X)$,
i.e., $\mu \in \Phi_A$.
Therefore $\Phi_{A+B} \subseteq \Phi_A$.
Conversely,
if $\mu \in \Phi_A$,
then the operator (20) is in $\Psi(\tilde{X},X)$.
Hence the same is true of $A+B- \mu$,
i.e.,
$\mu \in \Phi_{A+B}$.
Thus $\Phi_A = \Phi_{A+B}$.
This proves (15).
If $i(A+B-\lambda) = 0$,
w have by (l8)
\[
i(A+B-\mu) = i(A-\mu)
\]
and hence (l6) holds.

\begin{theorem1} % theorem 15
If $B$ is $A^2$-compact and $\lambda \in \Phi_A \cap \Phi_{A+B}$,
then (15) holds.
If, in addition, $i(A-\lambda) \eq i(A+B-\lambda)$, then (16) holds.
\end{theorem1}

Proof.
Consider the norm
\[
\norm{x}_1 = \norm{x} + \norm{Ax} + \norm{Bx} + \norm{A^2x} + \norm{BAx}
\]
on $\domain{A^2}$.
We claim that $\domain{A^2}$ equipped with this norm
becomes a Banach space $\tilde{X}$.
In fact if
$x_n \to x$,
$Ax_n \to y$,
$Bx_n \to z$,
$A^2x_n \to u$,
$BAx_n \to v$,
then $x \in \domain{A}$ since $A$ is closed
and $Ax = y$,
while $y \in \domain{A}$ and $Ay = u$ for the same reason.
Thus $x \in \domain{A^2}$.
Since $A+B$ is closed,
$(A+B)x = y+z$ and hence $Bz = z$.
The same reasoning gives $(A+B)Ax = u + v$,
whence $BAx = v$.
Now suppose $\mu \in \Phi_A$.
Then the operator $(A + B - \lambda) \in \Psi(X,X)$ by lemma 1
and hence is in $\Psi(\tilde{X},X)$.
However,
when we consider $B$ as an operator from $\tilde{X}$ to $X$,
it becomes a compact operator.
Since
\begin{equation}\tag{21}
(A+B-\lambda)(A-\mu) - (A+B-\mu)(A-\lambda) = (\lambda-\mu)B,
\end{equation}
we see that $(A + B - \mu)(A - \lambda)$ is a Fredholm operator
from $\tilde{X}$ to $X$ (theorem 1).
Thus it is in $\Psi(X,X)$ (lemma 2).
Since
$(A-\lambda) \in \Psi(X,X)$ we can apply lemma 4 to obtain that
$(A + B - \mu) \in \Psi(X,X)$.
Hence $\Phi_A \subset \Phi_{A+B}$.
The same reasoning applied in the opposite direction gives
$\Phi_A = \Phi_{A+B}$.
Moreover,
for $\mu \in \Phi_A$ we have by (21), lemma 1 and theorem 1
\[
i(A+B-\lambda) + i(A-\mu) = i(A+B-\mu) + i(A-\lambda).
\]
Hence
\[
i(A-\mu) = i(A+B-\mu).
\]
Thus $\sigma_{em}(A + B) \eq \sigma_{em}(A)$ and the proof is complete.

\paragraph*{}
We see from the proof of theorem 15 that the assumption
that $B$ is $A$-compact can be considerably weakened.

%------------------------------------------------------------------------------%
\section{Remarks}                                                              %
%------------------------------------------------------------------------------%
\begin{itemize}
\item[(R-1)]
In the Russian literature operators satisfying
properties 1, 3-5 are called $\Phi$-operators, with the $\Phi$ standing
for Fredholm (cf. [4]). The term Fredholm operator seems
to be reserved for $\Phi$-operators having index 0.
We have added assumption 2 for convenience and do not differentiate
on the basis of index.

\item[(R-2)]
Observations (O-4) and (O-5) are due to Atkinson [1] for bounded operators
(cf. also [11,12]).

\item[(R-3)]
Lemma 1 is also due to Atkinson [1] for bounded operators.
For unbounded operators it is due to Gohberg [3])
although there seems to be a gap in his proof (cf. A.M.S. translation of [4]).
The first part of our proof follows that of [4],
but our proof that $\range{AB}$ is closed is different.
Our proof that $AB$ is a closed operator was taken from Kato [8].
Other proofs may be found in [6,8].

\item[(R-4)]
For the histories of theorems 1 and 2 see [4],
Our proofs borrow ideas from Seeley [12] .
Our proof of $\alpha(A+T) < \alpha(A)$ appears to be much simpler
than any found in the literature.

\item[(R-5)]
The idea for lemma 2 and its proof came from Kato [8].

\item[(R-6)]
Theorems 3 and 4 as well as the device employed in
obtaining them from theorems 1 and 2 are due to Nagy [9].

\item[(R-7)]
Generalizations of theorems 5-7 can be found in [8,7].
Our proof of theorem 7 is taken directly from [4].

\item[(R-8)]
Lemmas 4 and 5 are apparently new.

\item[(R-9)]
The term essential spectrum originated in [5] where
it was applied to self-adjoint problems for ordinary differential equations
on a half-line.
In this case the essential spectrum is that part of the spectrum
which remains invariant under changes in the boundary conditions.
Browder [2] has proposed still another definition.

\item[(R-10)]
Theorems 12, 14, 15 are apparently new.
Similar results my be found in [2,14].
\end{itemize}

%------------------------------------------------------------------------------%
%                                 BIBLIOGRAPHY                                 %
%------------------------------------------------------------------------------%
\newpage
\fancyhead[C] {}
\fancyhead[OL]{}
\fancyhead[ER]{}
\thispagestyle{empty}

\titleformat{\section}
{\normalfont \bfseries \filcenter}
{\thesection.\;}{0mm}{}

\newlength{\bibitemsep}\setlength{\bibitemsep}{.5\baselineskip plus .05\baselineskip minus .05\baselineskip}
\newlength{\bibparskip}\setlength{\bibparskip}{0pt}
\let\oldthebibliography\thebibliography
\renewcommand\thebibliography[1]{%
  \oldthebibliography{#1}%
  \setlength{\parskip}{\bibitemsep}%
  \setlength{\itemsep}{\bibparskip}%
}

\vspace{8mm}
\begin{thebibliography}{99}
%\setlength\bibitemsep{3\itemsep}
\begin{small}

\bibitem{1}
F.V. Atkinson, Mat. Sb. 28(70) (1951), pp.\,3-14,

\bibitem{2}
F.E. Browder, Math, Ann. 142 (1961) pp.\,22-130.

\bibitem{3}
I.C. Gohberg, Mat. Sb. 53(75) (1953) pp.\,193-198.

\bibitem{4}
I.C. Gohberg and M.G. Krein, Usp. Mat. Nauk 12 (1957) pp.\,43-118;
Transl. A.M.S., 13(1960) p. 185-204 .

\bibitem{5}
Philip Hartman and Aurel Wintner, Amer. J. Math. 72 (1950) pp.\,543-552.

\bibitem{6}
J. T. Joichi, dissertation, Illinois, 1959.

\bibitem{7}
Shmuel Kaniel and Martin Schechter, Comm. Pure Appl.  Math., 16 (1965) pp.\,423-448.

\bibitem{8}
Tosio Kato, J. Analyse Math., 6 (1958) pp.\,261-322.

\bibitem{9}
B\'{e}la Sz.-Nagy, Acta Sci. Math. Szeged 14 (1951) pp.\,125-137.

\bibitem{10}
Martin Schechter, Invariance of the essential spectrum, to appear.

\bibitem{11}
J.T. Schwartz, Comm. Pure Appl. Math., 15 (1962), pp.\,75-90.

\bibitem{12}
R.T. Seeley, J. Math. Ann. Appl., 7 (1965) pp.\,289-309.

\bibitem{13}
Hermann Weyl, Rend. Circ. Mat. Palermo, 27 (1909) pp.\,373-392.

\end{small}
\end{thebibliography}
%------------------------------------------------------------------------------%
%                               END OF DOCUMENT                                %
%------------------------------------------------------------------------------%
\end{document}

NEW YORK UNIVERSITY
COURANT INSTITUTE - LIBkARY
IMM-NYU 335
JANUARY, 1965
NEW YORK UNIVERSITY
COURANT INSTITUTE OF
MATHEMATICAL SCIENCES
FREDHOLM OPERATORS AND THE
ESSENTIAL SPECTRUM
MARTIN SCHECHTER
PREPARED UNDER
CONTRACT NO. DA31-124-AR0- D-204
WITH THE
U. S. ARMY RESEARCH OFFICE

TI^™I-NYU 355
Courant Institute of Mathematical Sciences
Nev; York Un5.verslty
PREDHOLM OPERATORS AND THE
ESSENTIAL SPECTRUI4
Contract No. DA31-12^-AR0".D-204
January 1955
Martin Schechter
Department of Army Project No. 20011501B70^
"Request for additional copies by Agencies of the Department
of Defense, their contractors, and other Government agencies
should be directed to:
Defense Documentation Center
Cameron Station
Alexandria, Virginia 2231H
Department of Defense contractors must have established for
DDC services or have their "need-to-know" certified by the
cognizant military agencies of their project or contract."
Reproduction in I'jhole or in part is permitted for any purpose
of the United States Government.
